% =============================================================================
% Module 4: Solutions to Exercises
% =============================================================================
% This file is conditionally included based on \ifsolutions
% =============================================================================

\section{Solutions to Exercises}
\label{sec:lens_eq_solutions}

\noindent
Full worked solutions are also available in the Mathematica script
\solnlink{04_Lens_Equation/problems\_04.wl}{problems\_04.wl},
which verifies each answer symbolically and numerically.

% ----- Exercise 4.1 -----
\subsection*{Exercise~\ref{ex:derive_lens_eq}: Derivation of the Lens Equation}

\textbf{(a)} From the lensing geometry (Figure~\ref{fig:lensing_geometry}),
consider a light ray that passes through the lens plane at physical position
$\vecxi$ and is deflected by the angle $\boldsymbol{\hat{\alpha}}$.  In the
small-angle approximation, the physical position in the source plane is
determined by two contributions: (i) the extrapolation of the undeflected
ray from the lens plane to the source plane, which gives a position
$(\Ds/\Dd)\,\vecxi$, and (ii) the deflection, which displaces the ray by
$-\Dds\,\boldsymbol{\hat{\alpha}}(\vecxi)$ in the source plane.  Therefore:
\begin{equation*}
    \veceta = \frac{\Ds}{\Dd}\,\vecxi
    - \Dds\,\boldsymbol{\hat{\alpha}}(\vecxi) \,.
\end{equation*}

\textbf{(b)} Substituting $\vecxi = \Dd\,\vectheta$ and
$\veceta = \Ds\,\vecbeta$:
\begin{equation*}
    \Ds\,\vecbeta = \frac{\Ds}{\Dd}\,\Dd\,\vectheta
    - \Dds\,\boldsymbol{\hat{\alpha}}(\Dd\,\vectheta)
    = \Ds\,\vectheta - \Dds\,\boldsymbol{\hat{\alpha}}(\Dd\,\vectheta) \,.
\end{equation*}
Dividing by $\Ds$:
\begin{equation*}
    \vecbeta = \vectheta
    - \frac{\Dds}{\Ds}\,\boldsymbol{\hat{\alpha}}(\Dd\,\vectheta) \,.
\end{equation*}

\textbf{(c)} Defining the reduced deflection angle
$\vecalpha(\vectheta) \equiv (\Dds/\Ds)\,
\boldsymbol{\hat{\alpha}}(\Dd\,\vectheta)$, the lens equation becomes:
\begin{equation*}
    \vecbeta = \vectheta - \vecalpha(\vectheta) \,. \qquad \square
\end{equation*}

% ----- Exercise 4.2 -----
\subsection*{Exercise~\ref{ex:point_mass_solve}: Point Mass Lens Equation}

\textbf{(a)} For a point mass $M$ at the origin, the physical deflection
angle is $\alphahat = 4\Gnewton M / (\speedoflight^2 \xi)$ where
$\xi = \Dd\,\theta$.  The reduced deflection angle is:
\begin{equation*}
    \alpha(\theta) = \frac{\Dds}{\Ds}\,
    \frac{4\Gnewton M}{\speedoflight^2\,\Dd\,\theta}
    = \frac{1}{\theta}\,
    \frac{4\Gnewton M}{\speedoflight^2}\,\frac{\Dds}{\Dd\,\Ds}
    = \frac{\thetaE^2}{\theta} \,.
\end{equation*}
The lens equation is:
$\beta = \theta - \alpha(\theta) = \theta - \thetaE^2/\theta$.

\textbf{(b)} Multiplying by $\theta$:
$\theta^2 - \beta\,\theta - \thetaE^2 = 0$.
By the quadratic formula:
\begin{equation*}
    \theta_\pm = \frac{\beta \pm \sqrt{\beta^2 + 4\,\thetaE^2}}{2} \,.
\end{equation*}

\textbf{(c)} The discriminant is
$\Delta = \beta^2 + 4\,\thetaE^2 > 0$ for all $\beta$ (since $\thetaE > 0$).
Therefore two real solutions always exist.

\textbf{(d)} From Vieta's formulas for $\theta^2 - \beta\,\theta - \thetaE^2 = 0$:
\begin{itemize}[nosep]
    \item Sum of roots: $\theta_+ + \theta_- = \beta$. \checkmark
    \item Product of roots: $\theta_+\,\theta_- = -\thetaE^2$. \checkmark
\end{itemize}

\textbf{(e)} For $\beta > 0$: since $\sqrt{\beta^2 + 4\thetaE^2} > \beta$,
we have $\theta_+ > 0$.  Moreover,
$\theta_+ = (\beta + \sqrt{\beta^2 + 4\thetaE^2})/2 > \sqrt{\thetaE^2} = \thetaE$
because $\beta + \sqrt{\beta^2 + 4\thetaE^2} > 2\thetaE$ (square both
sides to verify).

For $\theta_-$: $\sqrt{\beta^2 + 4\thetaE^2} > \beta$ implies
$\theta_- = (\beta - \sqrt{\beta^2 + 4\thetaE^2})/2 < 0$.
Also $|\theta_-| = (\sqrt{\beta^2 + 4\thetaE^2} - \beta)/2 < \thetaE$,
which can be verified by noting $\sqrt{\beta^2 + 4\thetaE^2} - \beta < 2\thetaE$.
\hfill $\square$

% ----- Exercise 4.3 -----
\subsection*{Exercise~\ref{ex:einstein_radius_compute}: Einstein Radius for Various Systems}

Using $\thetaE = \sqrt{(4\Gnewton M/\speedoflight^2)\,(\Dds/\Dd\Ds)}$
with $\Gnewton = 6.674 \times 10^{-11}$~m$^3$\,kg$^{-1}$\,s$^{-2}$,
$\speedoflight = 2.998 \times 10^{8}$~m\,s$^{-1}$:

\medskip
\noindent\textbf{(a) Stellar microlensing:}
$M = 1.989 \times 10^{30}$~kg, $\Dd = \Dds = 5$~kpc $= 1.543 \times 10^{20}$~m,
$\Ds = 10$~kpc $= 3.086 \times 10^{20}$~m.
\begin{align*}
    \thetaE &= \sqrt{\frac{4 \times 6.674 \times 10^{-11} \times
    1.989 \times 10^{30}}{(2.998 \times 10^8)^2} \times
    \frac{1.543 \times 10^{20}}{1.543 \times 10^{20} \times
    3.086 \times 10^{20}}} \\
    &= \sqrt{5.906 \times 10^{-10} \times
    \frac{1}{3.086 \times 10^{20}}} \\
    &= \sqrt{1.914 \times 10^{-30}} \\
    &\approx 1.38 \times 10^{-15}~\text{rad}
    \approx 0.28~\text{mas} \,.
\end{align*}
Physical Einstein radius: $\RE = 5 \times 3.086 \times 10^{19} \times
1.38 \times 10^{-15} \approx 2.1 \times 10^{-10}$~Mpc
$\approx 6.6 \times 10^{-5}$~pc $\approx 14$~AU.

\medskip
\noindent\textbf{(b) Galaxy lensing:}
$M = 10^{12} \times 1.989 \times 10^{30}$~kg, $\Dd = 900$~Mpc,
$\Ds = 1700$~Mpc, $\Dds = 1200$~Mpc.
Converting to meters ($1~\text{Mpc} = 3.086 \times 10^{22}$~m):
\begin{align*}
    \thetaE &= \sqrt{\frac{4 \times 6.674 \times 10^{-11} \times
    1.989 \times 10^{42}}{(2.998 \times 10^8)^2} \times
    \frac{1200}{900 \times 1700}~\text{Mpc}^{-1}} \\
    &\approx 5.3 \times 10^{-6}~\text{rad} \approx 1.1'' \,.
\end{align*}
Physical Einstein radius: $\RE = 900 \times 5.3 \times 10^{-6}~\text{Mpc}
\approx 4.8 \times 10^{-3}$~Mpc $\approx 4.8$~kpc.

\medskip
\noindent\textbf{(c) Cluster lensing:}
$M = 10^{15} \times 1.989 \times 10^{30}$~kg, $\Dd = 900$~Mpc,
$\Ds = 1800$~Mpc, $\Dds = 1400$~Mpc.
\begin{equation*}
    \thetaE \approx 1.7 \times 10^{-4}~\text{rad} \approx 35'' \,.
\end{equation*}
Physical Einstein radius: $\RE \approx 150$~kpc.

% ----- Exercise 4.4 -----
\subsection*{Exercise~\ref{ex:sigma_cr_compute}: Critical Surface Mass Density}

\textbf{(a)} Dimensional analysis of $\Sigmacr = \speedoflight^2 /
(4\pi\Gnewton) \times \Ds / (\Dd\,\Dds)$:
\begin{equation*}
    [\Sigmacr] = \frac{(\text{m/s})^2}{(\text{m}^3\,\text{kg}^{-1}\,\text{s}^{-2})}
    \times \frac{\text{m}}{\text{m} \times \text{m}}
    = \frac{\text{m}^2\,\text{s}^{-2} \times \text{kg}\,\text{s}^{2}}
    {\text{m}^3} \times \frac{1}{\text{m}}
    = \frac{\text{kg}}{\text{m}^2} \,.
\end{equation*}
So $\Sigmacr$ has units of mass per area, as required. \checkmark

\textbf{(b)} For the galaxy lens ($\Dd = 900$~Mpc, $\Ds = 1700$~Mpc,
$\Dds = 1200$~Mpc):
\begin{equation*}
    \Sigmacr = \frac{(2.998 \times 10^{8})^2}{4\pi \times
    6.674 \times 10^{-11}} \times
    \frac{1700 \times 3.086 \times 10^{22}}{900 \times 3.086 \times 10^{22}
    \times 1200 \times 3.086 \times 10^{22}}
    \approx 3.5~\text{kg\,m}^{-2} \,.
\end{equation*}
Converting: $3.5~\text{kg\,m}^{-2} = 0.35~\text{g\,cm}^{-2}
\approx 1700~\Msun\,\text{pc}^{-2}$.

\textbf{(c)} For the cluster lens ($\Dd = 900$~Mpc, $\Ds = 1800$~Mpc,
$\Dds = 1400$~Mpc):
\begin{equation*}
    \Sigmacr \approx 3.0~\text{kg\,m}^{-2} = 0.30~\text{g\,cm}^{-2}
    \approx 1500~\Msun\,\text{pc}^{-2} \,.
\end{equation*}

\textbf{(d)} A typical elliptical galaxy with
$\Sigma_0 \approx 5000~\Msun\,\text{pc}^{-2}$ has
$\kappab = \Sigma_0/\Sigmacr \approx 5000/1700 \approx 3 > 1$.
Yes, it is a strong lens --- the central surface mass density exceeds
the critical value, enabling the formation of multiple images.

% ----- Exercise 4.5 -----
\subsection*{Exercise~\ref{ex:image_positions_plot}: Image Positions for a Point Mass Lens}

\textbf{(a)} For $\thetaE = 1''$ and $\theta_\pm = (\beta \pm
\sqrt{\beta^2 + 4})/2$ (in arcseconds):

\begin{center}
\begin{tabular}{rccccc}
\toprule
$\beta$ (arcsec) & $\theta_+$ (arcsec) & $\theta_-$ (arcsec)
& $\Delta\theta$ (arcsec) & $|\mu_+|$ & $|\mu_-|$ \\
\midrule
0.1  & $+1.025$ & $-0.975$ & 2.000 & 5.13 & 4.63 \\
0.5  & $+1.281$ & $-0.781$ & 2.062 & 1.64 & 0.64 \\
1.0  & $+1.618$ & $-0.618$ & 2.236 & 1.17 & 0.17 \\
2.0  & $+2.414$ & $-0.414$ & 2.828 & 1.06 & 0.06 \\
5.0  & $+5.193$ & $-0.193$ & 5.385 & 1.01 & 0.01 \\
\bottomrule
\end{tabular}
\end{center}

\textbf{(b)--(d)} Plots are generated in
\mathlink{04_Lens_Equation/lens_equation.wl}{lens\_equation.wl}
and
\solnlink{04_Lens_Equation/problems\_04.wl}{problems\_04.wl}.
See Figures~\ref{fig:image_positions} and~\ref{fig:einstein_ring}.

The minimum image separation is $\Delta\theta_{\mathrm{min}} = 2\thetaE = 2''$,
occurring at $\beta = 0$ (Einstein ring).

The total magnification $|\mu_+| + |\mu_-|$ diverges as $\beta \to 0$
and approaches 1 (no magnification) as $\beta \to \infty$.  For
$\beta = \thetaE$, the total magnification is
$(1 + 2)/(1 \cdot \sqrt{5}) \approx 1.34$.
